% Replaces main.tex Section 3.1. Reserve Proposition 1 for Appendix A.
\subsection{Utility mismatch under ratio risk}
Let $X$ be the information available when acceptance is chosen. For
nonnegative integrable loss $L$ and exposure $W$, define
\[
 \ell(X)=\E[L\mid X],\quad w(X)=\E[W\mid X],\quad T=\E W>0,
 \quad a(X)\in[0,1].
\]
Row coverage, exposure coverage, and accepted ratio risk are
\begin{equation}
 c(a)=\E a,\qquad d_W(a)=\frac{\E[aW]}{T},\qquad
 R(a)=\frac{\E[aL]}{\E[aW]}.
 \label{eq:ratio}
\end{equation}
The accepted denominator must be positive. The contract is $R(a)\leq r$
and $c(a)\geq c_0$. Its signed excess is
\begin{equation}
 g_r=L-rW,\qquad m_r=\ell-rw,\qquad
 R(a)\leq r\ \Longleftrightarrow\ \E[a m_r]\leq0.
 \label{eq:excess}
\end{equation}
WAPE is the specialization $L=|Y-f(X)|$, $W=Y$, with $Y,f\geq0$;
then $\ell=e$, $w=\mu$, and $d_W=d$ is demand coverage.

Fixed-row-coverage excess is minimized by accepting the smallest
$m_r$ values, with randomization on ties (Proposition~\ref{prop:threshold}
in Appendix~\ref{app:proofs}). To maximize exposure instead, change measure
to $d\nu=w\,dP_X/T$ after rejecting $w=0,\ell>0$ contexts: then $d_W(a)=\E_\nu a$ and
$R(a)=\E_\nu[a\ell/w]/\E_\nu a$.
\setcounter{proposition}{1}
\begin{proposition}[Exposure-weighted ranking]
\label{prop:weighted}
If a positive-exposure feasible rule exists and the row floor is
nonbinding, exposure coverage is optimized by thresholding $\ell/w$,
with boundary randomization. Reject $w=0,\ell>0$ contexts; $w=\ell=0$
contexts are neutral for this objective. The ranking is independent of
$r$; its threshold is not.
\end{proposition}
Thus the population rankings are $\ell-rw$ for rows and $\ell/w$ for
exposure. A fixed ranking's largest feasible threshold maximizes either
coverage because its acceptance sets are nested; objective comparisons
must therefore also compare ranking families (Appendix~\ref{app:objective}).

For disjoint context sets $K,U$, write
\[
 M_A=\E[\ind_Aw],\qquad B=-\E[\ind_Km_r],\qquad
 M_K>0,\quad B\geq0.
\]
Here $B$ is the excess slack supplied by a common accepted set.
\begin{proposition}[Sharp objective reversal for ratio risk]
\label{prop:lowdemand}
Fix $r>0$, $\rho>r$, and $\varepsilon>0$. If $\ell=\rho w$ on $U$ and
$0<M_U\leq\varepsilon\Pr(U)$, then $R(\ind_U)=\rho$ and
$K\cup U$ satisfies the cap exactly when $(\rho-r)M_U\leq B$.
Accepting $U$ adds $\Pr(U)$ row coverage but at most
$\varepsilon\Pr(U)/T$ exposure coverage.

For every $0\leq c_0<1$ and $\eta>0$, there exists a three-context
law with $w>0$, $0\leq\ell\leq\rho w$, and exact conditional moments
whose unique row- and exposure-optimal policies $a_R,a_W$ satisfy
the cap and the nonbinding row floor, and obey
\[
 c(a_R)-c(a_W)>1-c_0-\eta,\qquad
 d_W(a_W)-d_W(a_R)>1-\eta.
\]
The universal bounds $1-c_0$ and $1$ are jointly sharp. For
$0<r<1$ and $\rho=1$, this sharpness also holds with binary $0\leq L\leq W\leq1$.
\end{proposition}
The sharpness statement ranges over joint laws that can realize the
conditional moments in Appendix~\ref{app:proofs}; it is not a claim
about every fixed metric class. If $W$ is constant, for example,
$d_W=c$ and reversal is impossible. For a fixed multi-label classifier, whole-example abstention yields an exact micro-precision reversal (Appendix~\ref{app:classification}). Ordinary binary precision with an unrestricted gate instead has a common optimum for both coverage objectives. For whole-example micro-precision with $1\leq w\leq m$, each coverage gap is at most $(m-1)/(m+1)$; Appendix~\ref{app:bounded_exposure} sharpens this to a joint bound.

\paragraph{WAPE corollary and estimation.}
For $0<r<1$, setting $\rho=1$, $W=Y$, and $f=0$ on $U$ gives
$e=\mu$: near-zero-demand rows have budget cost $(1-r)\mu$ yet
unit row reward. The three-context law is realizable by deterministic
nonnegative forecasts (Appendix~\ref{app:proofs}), so both sharp bounds
hold for WAPE. A low-exposure share alone cannot bound the forgone
exposure without a positive excess-rate margin; the appendix gives
that bound and the exact budget identity. More generally $|e-\mu|\leq f$.
Exact-head optimality does not imply estimated-ratio robustness.
Section~\ref{sec:objectives} and Appendix~\ref{app:proofs} distinguish
objective alignment from small-denominator error amplification.
